\documentclass[11pt]{article}
\usepackage[margin=1in]{geometry}
\usepackage{booktabs}
\usepackage{enumitem}
\usepackage{hyperref}
\usepackage{microtype}
\usepackage{amsmath}

\title{Lab Notebook --- Definitive Anchoring Experiment, Stage A}
\author{Alessa Carbo}
\date{2026-07-28}

\begin{document}
\maketitle

\section{Objective and scope}

This entry records Stage A of the definitive anchoring experiment, which replaces
the incremental pilot sequence recorded in week 8. The experiment installs one
counterfactual set through three matched corpora---premise-position (P),
assertion (A), and entity-decoupled (E)---and measures collateral drift against
graph distance (H1) and recovery cost under a $\lambda$-swept adaptive relearning
attacker (H2).

Stage A has two jobs: a retrospective check of the selection heuristic on the
pilot's 15 claims (A.5), and a ranked candidate table from which the final target
set is chosen by hand (A.2--A.4). Three amendments were adopted before running:

\begin{enumerate}[nosep]
  \item[\textbf{A1}] The drift-versus-distance curve stays within a single
        instrument across \texttt{crossing}, \texttt{near-retain}, and
        \texttt{far-retain}. \texttt{off-domain-retain} is a different
        instrument and is reported separately, never as a point on the curve.
  \item[\textbf{A2}] The Stage-H calibration gains a third control: a
        fabricated target, never true and never installed, giving a
        false-positive floor on measured ``recovery.''
  \item[\textbf{A3}] A null in A.5 is uninformative at $n=15$, not
        disconfirming. A clear positive raises confidence in the ranker; a null
        means proceed with the ranker while treating
        support-as-installability as an untested assumption to be checked at
        $n\approx40$--$60$, where there is power.
\end{enumerate}

\section{Graph convention}

Verified before any analysis, because every Stage-A quantity depends on it. In
\texttt{viral\_vectors\_graph.json}, edges run
\[
\text{node} \longrightarrow \text{its prerequisite},
\]
so \texttt{in\_degree}$(n)$ counts $n$'s immediate \emph{dependents}, and the
claims that depend on $n$ transitively are $\mathrm{ancestors}(n)$, not
descendants. The brief's ``downstream descendant closure'' is therefore an
ancestor set under this orientation. Depth of a dependent is its shortest path
length to the root. Confirmed against claim text: the top hubs by dependent count
are C375 (51), C009 (44), C031 (41), C049 (37), C034 (26)---none of which is in
the pilot's edited set.

\section{Graph structure relevant to selection}

\begin{table}[h]
\centering
\begin{tabular}{lr}
\toprule
Quantity & Value \\
\midrule
Nodes & 533 \\
Edges & 1344 \\
Strongly connected components ($>1$ node) & 68 \\
Nodes inside an SCC & 230 (43.2\%) \\
Edges internal to an SCC & 429 / 1344 \\
SCC condensation & 371 units, acyclic \\
Largest undirected component & 508 / 533 \\
Remaining components & 4, 3, 3, 3, 2, 2, 2 \\
\bottomrule
\end{tabular}
\caption{The SCCs are mostly small (36 of size 2, 18 of size 3) but include one
of size 25 and singles of size 11, 10, 9, 7, and 6. Almost the entire graph is
one undirected component.}
\end{table}

\section{Stage A.5 --- retrospective check on the prior 15}

For each pilot-edited claim we computed dependent closure size, immediate
dependents, maximum and mean dependent depth, the number of dependents at depth
$\geq 2$, and total distinct dependent-to-root path count (simple paths, cutoff
6, per-dependent cap 1000). These were correlated against paired
$\Delta P(\mathrm{twin})$ at the selected six-epoch no-replay condition.

\begin{table}[h]
\centering
\begin{tabular}{lrrrrrr}
\toprule
Claim & Closure & Immediate & $d_{\max}$ & $n(d\geq2)$ & Paths & $\Delta P_t$ (e6) \\
\midrule
C005 & 1  & 1 & 1 & 0  & 1  & $+0.294$ \\
C017 & 1  & 1 & 1 & 0  & 1  & $+0.718$ \\
C028 & 0  & 0 & 0 & 0  & 0  & $+0.438$ \\
C065 & 5  & 1 & 3 & 4  & 5  & $+0.473$ \\
C072 & 2  & 2 & 1 & 0  & 3  & $+0.426$ \\
C097 & 33 & 2 & 7 & 31 & 42 & $+0.187$ \\
C101 & 0  & 0 & 0 & 0  & 0  & $-0.246$ \\
C126 & 0  & 0 & 0 & 0  & 0  & $-0.319$ \\
C178 & 2  & 2 & 1 & 0  & 2  & $+0.384$ \\
C305 & 0  & 0 & 0 & 0  & 0  & $-0.266$ \\
C353 & 4  & 1 & 2 & 3  & 4  & $+0.002$ \\
C387 & 0  & 0 & 0 & 0  & 0  & $-0.242$ \\
C388 & 0  & 0 & 0 & 0  & 0  & $+0.138$ \\
C419 & 0  & 0 & 0 & 0  & 0  & $+0.607$ \\
C476 & 0  & 0 & 0 & 0  & 0  & $+0.304$ \\
\bottomrule
\end{tabular}
\caption{Support structure of the pilot's edited set. Eight of 15 claims have no
downstream dependents whatsoever; only C065, C097, and C353 have a dependent at
depth $\geq2$.}
\end{table}

\subsection{Prespecified correlation: null}

\begin{table}[h]
\centering
\begin{tabular}{lrrrr}
\toprule
Predictor & $\rho$ (e6) & perm $p$ & $\rho$ (e3) & $\rho$ (e12) \\
\midrule
Closure size          & $+0.316$ & 0.242 & $-0.004$ & $+0.214$ \\
Immediate dependents  & $+0.362$ & 0.181 & $-0.004$ & $+0.291$ \\
Max dependent depth   & $+0.327$ & 0.225 & $-0.031$ & $+0.192$ \\
$n$ at depth $\geq2$  & $+0.049$ & 0.872 & $-0.112$ & $-0.097$ \\
Total path count      & $+0.318$ & 0.240 & $-0.021$ & $+0.204$ \\
\bottomrule
\end{tabular}
\caption{Spearman correlations with paired $\Delta P(\mathrm{twin})$;
permutation $p$ from 10{,}000 deterministic relabelings (seed 42). Kendall
$\tau_b$ agrees throughout. No scipy in this environment, so the rank
statistics, tie handling, and permutation test are implemented locally.}
\end{table}

Directionally positive at six and twelve epochs, absent at three, and
significant nowhere. Per amendment A3 this is treated as uninformative.

\subsection{Why it is uninformative: the predictor is degenerate}

The stronger finding is structural rather than statistical. The support variable
has almost no variation to correlate against: 8 of 15 claims sit at exactly zero,
the largest tie group is 53\% of the sample, and only 3 of 15 have any dependent
at depth $\geq 2$. A rank correlation on that is dominated by tie handling.

The consequence is sharper than a power complaint. \textbf{The pilot could not
have observed propagated coherence regardless of install strength}, because for
most of its edited claims there was nothing downstream to propagate into.
Acquisition in week 8 was measured on bare roots. This is exactly the floor
problem the Stage-A selection fix addresses, and it is independently motivated by
the graph structure rather than by any correlation result.

\subsection{Post-hoc observation (not prespecified)}

Inspecting the table, the informative contrast is not the ordering within a
53\%-tied variable but the binary any-support versus no-support split. Every
claim with a dependent moved toward the twin; every negative mover had none.

\begin{table}[h]
\centering
\begin{tabular}{llrrrr}
\toprule
Condition & Set & Support & No support & Perm $p$ & Fisher $p$ \\
\midrule
e6, all 15      & mean $\Delta P_t$ & $+0.355$ (7/7 pos) & $+0.052$ (4/8 pos) & 0.085 & 0.051 \\
e6, excl. C028  & mean $\Delta P_t$ & $+0.355$ (7/7 pos) & $-0.003$ (3/7 pos) & 0.050 & 0.035 \\
e12, all 15     & mean $\Delta P_t$ & $+0.233$ (6/7 pos) & $-0.005$ (4/8 pos) & 0.271 & 0.182 \\
e3, all 15      & mean $\Delta P_t$ & $+0.105$ (4/7 pos) & $+0.052$ (4/8 pos) & 0.740 & 0.595 \\
\bottomrule
\end{tabular}
\caption{Post-hoc split. Fisher $p$ is the one-sided probability that all
supported claims land in the positive group by chance.}
\end{table}

\textbf{This cut was chosen after seeing the data, carries no multiplicity
correction, and is one of several cuts that could have been taken.} It is
recorded as hypothesis-generating only, and it does not upgrade the prespecified
null. It does point the same direction the ranker predicts, and it is testable
with power at $n\approx40$--$60$.

\paragraph{Correction to the brief's motivation.}
The brief cites ``8/15 moved toward twin, 7/15 away.'' That is the three-epoch
run, which was superseded. At the selected six-epoch no-replay condition the
split is 11 positive / 4 negative. The heterogeneity argument stands but is
milder than stated.

\section{Stage A.2--A.4 --- ranked candidates and a specification conflict}

Candidate roots were profiled on all four brief columns, with eligibility
requiring a dependent at depth $\geq2$, cycle-cleanliness, and a
non-procedural claim. A distribution screen flags refusal-prone procedural text;
it is calibrated on C028 (\emph{``Add 15 uM of the inhibitor at the time of viral
transfection and every 4 hours\ldots''}), whose twin came back as a refusal. An
initial version of the screen also hard-flagged claims lacking a recognized
declarative marker, which false-positived on four ordinary factual claims
(C045, C046, C009, C091) whose main verbs---``couple,'' ``promote,'' ``present,''
``stabilize''---were simply outside the marker list. That test was demoted to a
soft note; the hard flag is now imperative-initial only, which catches C028
exactly.

\subsection{The conflict}

Two brief requirements are mutually unsatisfiable on this graph:
cycle-cleanliness (``the closure must not intersect the 68 SCCs'') and target
scale ($n\approx40$--$60$ claims from roots plus high-support dependents). With
43\% of nodes inside an SCC and 508 of 533 in one undirected component, every
root with a deep closure reaches a cycle, and every cycle-clean root sits on a
shallow periphery.

\begin{table}[h]
\centering
\begin{tabular}{lrrrl}
\toprule
Cycle policy & Eligible & Candidates & Max union & Target band 40--60 \\
\midrule
\texttt{strict}      & 11 & 533 & 14  & unreachable \\
\texttt{condense}, cap 1 & 26 & 371 & $\geq60$ & reached at 3 roots \\
\texttt{condense}, cap 2 & 35 & 371 & $\geq60$ & reached at 1 root \\
\bottomrule
\end{tabular}
\caption{Under the brief as written the largest achievable edited set is 14
claims, and the eligible roots have closures of 2--7 shallow dependents.}
\end{table}

\subsection{The two readings, both implemented}

\texttt{strict} is the brief as written. \texttt{condense} ranks on the SCC
condensation, which is acyclic, and requires only that every \emph{edited} unit
be an SCC no larger than \texttt{--max-scc-size}; the closure may traverse
condensed regions. Depth and support remain well defined either way, because the
condensation is a DAG. What \texttt{condense} relaxes is solely whether an
edited claim may be \emph{reached through} a cyclic region. At cap 1 everything
edited is SCC-free---identical to \texttt{strict} in what gets edited---and only
the closure constraint differs. Both rankings are persisted so the choice is
visible rather than buried.

\begin{table}[h]
\centering
\begin{tabular}{lrrrrrr}
\toprule
Root & Closure & Immediate & $d_{\max}$ & $n(d\geq2)$ & Paths/dep & Edit set \\
\midrule
C012 & 74 & 8  & 6 & 66 & 1.66 & 32 \\
C020 & 67 & 9  & 5 & 58 & 1.78 & 30 \\
C131 & 46 & 12 & 4 & 34 & 1.72 & 21 \\
C163 & 34 & 3  & 5 & 31 & 1.47 & 14 \\
C133 & 36 & 4  & 5 & 32 & 1.42 & 13 \\
C097 & 28 & 2  & 5 & 26 & 1.25 & 6  \\
\bottomrule
\end{tabular}
\caption{Top eligible roots under \texttt{condense} cap 1. C020's edit set is
nearly nested inside C012's (greedy union grows 32 $\to$ 33), so the two are near
substitutes rather than complements.}
\end{table}

\subsection{Cycle policy selected, with the cost instrumented}

\texttt{condense} at cap 1 was selected. Rather than assume the relaxation is
harmless, its specific cost is now measured. Cap 1's cost is a \emph{mid-chain
coherence gap}: a closure passing through a cyclic region propagates the
counterfactual through intermediate nodes that stayed real, and inside an SCC
``upstream'' is not well defined, so the gap cannot even be reasoned about
directionally. We therefore record, per edited node, whether its propagation path
from the selected root traverses an SCC, and prespecify the comparison in Stage G:
do traversal-crossing edited nodes install worse than clean-path ones? If yes, we
have a covariate to control for; if no, cycle traversal was harmless and we
learned it cheaply. Either way an unmeasured risk becomes a measured one.

Declaring this before install matters: it makes the Stage-G check prespecified
rather than post-hoc, which is precisely the weakness that limits what A.5's
support split can support.

The field is emitted per edited node with \texttt{n\_paths},
\texttt{n\_paths\_via\_scc}, \texttt{any\_path\_via\_scc},
\texttt{all\_paths\_via\_scc}, \texttt{frac\_paths\_via\_scc}, and the list of
SCC units on those paths. Only internal path nodes count---under cap 1 both
endpoints are SCC-free by construction, so an endpoint test would be vacuous. The
\emph{any} versus \emph{all} distinction carries the meaning: if some path is
clean a coherent chain exists, and only when every path traverses a cycle is the
gap unavoidable. Where a node is reachable from more than one selected root the
merge takes the most favourable record, since one coherent chain suffices.

\subsection{The covariate is near-constant in binary form}

Instrumenting it immediately showed that the binary version cannot carry an
inferential test on this graph.

\begin{table}[h]
\centering
\begin{tabular}{lrrr}
\toprule
Root set & Clean & Avoidable & Unavoidable \\
\midrule
\{C012, C131\}, 51 non-root edited nodes & 3 & 9 & 39 \\
13-root mixed set, 34 non-root edited nodes & 6 & 4 & 24 \\
\bottomrule
\end{tabular}
\caption{Binary SCC-traversal split. The cycle-free periphery of this graph
contains very few high-support dependents, so the clean group stays small however
the roots are chosen; a deliberately mixed root set built from all 11 cycle-clean
peripheral roots only moves it from 3 to 6.}
\end{table}

A 3-versus-48 comparison is the same underpowered shape as A.5, so a null would
be uninformative for the same reason. The graded form does have usable variance:

\begin{table}[h]
\centering
\begin{tabular}{lrrrrr}
\toprule
Distinct SCC units on propagation paths & 0 & 1 & 2 & 3 & 4 \\
\midrule
Edited nodes (\{C012, C131\}) & 3 & 5 & 13 & 20 & 10 \\
\bottomrule
\end{tabular}
\caption{\texttt{n\_scc\_units\_on\_paths} across the 51 non-root edited nodes.
A five-level ordinal with usable spread, versus a 3/48 binary split.
\texttt{frac\_paths\_via\_scc} is not usable (39 of 51 nodes sit at 1.0).}
\end{table}

Stage G therefore prespecifies the \textbf{graded} test as primary---rank-correlate
paired $\Delta P(\mathrm{twin})$ against the number of distinct cyclic regions the
propagation chain passes through, across all non-root edited nodes---and reports
the binary split descriptively alongside it, with per-group $n$. The graded test
uses every edited node and tests the same idea more finely: more cyclic regions
traversed should mean a larger coherence gap and a worse install.

Note that \{C012, C131\} is \emph{better} for this covariate than the mixed
13-root set (51 non-root nodes across five levels versus 34 more concentrated
ones), so no root-set contortion is needed to make the check work.

\subsection{Stage B strata preview for a viable pair}

\begin{table}[h]
\centering
\begin{tabular}{lr}
\toprule
Stratum & $n$ \\
\midrule
\texttt{forget} (edited)      & 53 \\
\texttt{crossing} (distance 1) & 70 \\
\texttt{near-retain} (2--3)   & 317 \\
\texttt{far-retain} ($\geq4$ or disconnected) & 93 \\
\bottomrule
\end{tabular}
\caption{Roots \{C012, C131\} under \texttt{condense} cap 1. Stratum distance is
undirected graph distance: for collateral-drift purposes a prerequisite of an
edited node is as adjacent as a dependent. \texttt{off-domain-retain} is a
separate instrument and is not enumerated here (amendment A1).}
\end{table}

This lands \texttt{forget} inside the 40--60 band with every graph stratum
populated well above a 30-claim floor, so the H1 distance curve would have
comparable counts at each point. \texttt{near-retain} is large enough to
subsample rather than use whole.

\section{Artifacts}

\begin{itemize}[nosep]
  \item \texttt{scripts/stage\_a\_step5\_support\_correlation.py} and
        \texttt{data/stage\_a\_step5\_support\_correlation.json}
  \item \texttt{scripts/stage\_a\_rank\_candidates.py}
  \item \texttt{data/stage\_a\_candidate\_ranking\_\{strict,condense1,condense2\}.json}
  \item \texttt{data/stage\_b\_strata\_preview\_condense1\_C012\_C131.json} and
        \texttt{\dots\_condense1\_mixed.json}, each carrying the per-edited-node
        SCC-traversal covariate
\end{itemize}

\section{Open decisions}

The brief's four, plus one surfaced by Stage A:

\begin{enumerate}[nosep]
  \item Final target node set (from the ranked table).
  \item Premise/assertion ratio within arm P.
  \item $\lambda$ grid for the adaptive attack.
  \item Attack-data provenance confirmation.
\end{enumerate}

Decision 1 remains open: the ranked table is the input to it, and no root set was
auto-selected. Decisions 2--4 are untouched.

\paragraph{Resolved during Stage A.}
A fifth decision surfaced---cycle policy---and was resolved to
\texttt{condense} at cap 1, on the condition that the relaxation's cost be
instrumented rather than assumed. \texttt{strict} was rejected because it caps the
edited set at 14 shallow claims, reintroducing a milder form of the floor problem
A.5 diagnosed; cap 2 was not needed, since cap 1 reaches the target band while
keeping every edited unit SCC-free.

\section{Limitations}

\begin{itemize}[nosep]
  \item A.5's prespecified test is null and low-powered; the post-hoc split is
        not evidence and is labelled as such.
  \item Path counts are simple paths with cutoff 6 and a per-dependent cap of
        1000, so \texttt{paths\_total} is a bounded support proxy, not an exact
        count.
  \item The distribution screen is a text heuristic. It catches C028's
        signature and no longer false-positives on the four claims checked, but
        it has not been validated against a labelled set, and every candidate in
        a final target set should be read by hand.
  \item Undirected stratum distance is a defensible default, not a validated
        choice; the directed variant would populate the strata differently.
  \item \texttt{off-domain-retain} remains unspecified. The MMLU-bio and
        WMDP-bio slices already wired up elsewhere in this repo are not
        off-domain with respect to a viral-vector graph, so a genuinely
        out-of-domain benchmark slice still needs to be named.
  \item The SCC-traversal covariate's binary form is near-constant (3 clean of 51
        edited nodes) and cannot support an inferential test on this graph
        whatever root set is chosen. The graded form is the powered version, but
        it is a proxy: number of cyclic regions crossed is not the same quantity
        as size of coherence gap, and the two could come apart.
  \item Path enumeration is capped, so a node at the 1000-path cap may have
        traversal statistics computed from a truncated path set. No node in the
        previewed sets is near the cap (maximum observed path count is 6), but
        the cap would bite on the large SCC-heavy roots excluded here.
\end{itemize}

\end{document}
